\documentclass[11pt]{article}

\usepackage[margin=1in]{geometry}
\usepackage{amsmath}
\usepackage{amssymb}
\usepackage{amsthm}
\usepackage{booktabs}
\usepackage{array}
\PassOptionsToPackage{hyphens}{url}
\usepackage[hidelinks]{hyperref}

\setlength{\emergencystretch}{3em}

\theoremstyle{plain}
\newtheorem{theorem}{Theorem}[section]
\newtheorem{lemma}[theorem]{Lemma}
\newtheorem{proposition}[theorem]{Proposition}
\newtheorem{corollary}[theorem]{Corollary}
\newtheorem{conjecture}{Conjecture}
\renewcommand{\theconjecture}{\arabic{conjecture}}

\theoremstyle{definition}
\newtheorem{definition}[theorem]{Definition}

\theoremstyle{remark}
\newtheorem{remark}[theorem]{Remark}

\newcommand{\R}{\mathbf{R}}
\newcommand{\Cx}{\mathbf{C}}
\newcommand{\N}{\mathbf{N}}
\newcommand{\re}{\operatorname{Re}}
\newcommand{\im}{\operatorname{Im}}
\newcommand{\tphi}{\widetilde{\varphi}}

\title{On Haglund's Conjecture~4 for the Riemann $\Xi$ approximants}
\author{Kunal Tyagi\thanks{%
\texttt{hello.jay.tyagi@gmail.com}. \emph{Use of AI.} The mathematics in this
paper --- the derivations, the computations, the verification suite, and the
text --- was produced by Claude (Anthropic) working under the author's
direction; the author set the objectives, made the technical and editorial
decisions, and is responsible for the content. Claude is a tool and is not an
author. The complete working record is public:
\url{https://github.com/x67ai/riemann-rh-program}.}}
\date{October 3, 2026}

\begin{document}

\maketitle

\begin{abstract}
Let $\Xi_k$ be the $k$-th partial sum of Riemann's expansion
$\Xi = \sum_{n \ge 1}\Phi_n$ of the $\Xi$-function in incomplete gamma
functions. Haglund conjectured (his Conjecture~4) that for every $k \ge 1$ the
imaginary part of each non-real zero of $\Xi_k + t\Phi_{k+1}$ decreases as $t$
goes from $0$ to $1$. We record what can be said about this pencil, each item
with its kind. Proved: the pencil is $\Xi - L_t$ with a level $L_t$ that is
positive on the real axis and decreases with $t$; its real zeros lie in the
positive lobes of $\Xi$, and real zeros stay real as $t$ increases if and only
if every critical point of $\Xi_{k+1}/\Phi_{k+1}$ on the real axis with value
in $(0,1]$ is a non-degenerate maximum; a simple zero $\beta$ of $\Xi$ in the
upper half-plane with $\im \Xi'(\beta) > 0$ would give, for every large $k$, a
zero of the pencil whose imaginary part strictly increases on $[0,1]$. Known in
substance (Csordas--Smith): if the zeros of $\Xi$ are real and simple and the
level is frozen at a constant, every non-real zero descends along its own curve
and lands at the top of its own lobe. Numerical, in floating point and not in
interval arithmetic, with two independent programs on $k = 1, 2, 3$: along every
branch followed in complete windows for $k \le 7$ and in windows around the
landing region for $k = 15, 20, 26, 27$, the imaginary part decreased at every
step, and the real-axis criterion holds for $k \le 50$ on $[0, 4(k+2)^2 + 40]$.
Heuristic: why the true, slowly varying level leaves the frozen picture intact.
\end{abstract}

%% SECTION 1
\section{The conjecture and what is shown here}\label{sec:intro}

\subsection{The pencil}\label{sec:pencil}
Write $s = \tfrac12 + iz$ and let
$\Xi(z) = \tfrac12 s(s-1)\pi^{-s/2}\Gamma(s/2)\zeta(s)$, which is even, entire
and real on the real axis. Riemann's series $\Xi = \sum_{n\ge1}\Phi_n$ and
Haglund's approximants $\Xi_N := \sum_{n \le N}\Phi_n$ are defined in
\cite[(1)--(14), pp.~1--3]{Hag} (page and equation numbers refer to the arXiv
version); see also \cite[(3)]{Tya}. From Haglund's definition (his (6) with
$z \in \Cx$, and his (14)), and on the real axis also from
\cite[Theorem~3.3 and its proof]{Tya},
\begin{equation}\label{eq:kernel}
\Phi_n(z) = 2\int_0^\infty \tphi_n(v)\cos(zv)\,dv \quad (z \in \Cx), \qquad
\tphi_n(v) = 2y(2y-3)\,e^{v/2 - y}, \quad y = \pi n^2 e^{2v},
\end{equation}
and for $n \ge 2$ the kernel $\tphi_n$ is positive, strictly decreasing and
strictly convex on $[0,\infty)$ \cite[proof of Theorem~3.3]{Tya}; for $n = 1$
it is positive but not monotone at $0$. Put
$Q_N := \sum_{n > N}\Phi_n = \Xi - \Xi_N$, an entire function (the series
converges locally uniformly \cite[(12)]{Hag}). For $k \ge 1$ the pencil is
\[
F_t := \Xi_k + t\,\Phi_{k+1}, \qquad 0 \le t \le 1, \qquad u := 1 - t ,
\]
so that $F_0 = \Xi_k$ and $F_1 = \Xi_{k+1}$. Haglund writes \cite[p.~11]{Hag}:
``Another interesting phenomena occurs when we consider the zeros of
$\Xi_k(z) + t\Phi_{k+1}(z)$. As we let $t$ vary continuously from $0$ to $1$,
computations indicate that the imaginary parts of the non-real zeros decrease
monotonically (i.e.\ continuously), in a very regular manner, until, for high
enough $k$, they collide with the corresponding zero (from Schwartz reflection)
in the fourth quadrant, and arrive on the real line, where they remain.'' He then
states:

\setcounter{conjecture}{3}
\begin{conjecture}[Haglund {\cite[p.~11]{Hag}}]\label{conj:4}
``For $k \ge 1$, the imaginary part of each non-real zero of
$\Xi_k(z) + t\Phi_{k+1}(z)$ decreases monotonically (i.e.\ continuously) as $t$
goes from $0$ to $1$.''
\end{conjecture}

The journal version \cite[p.~310]{HagJ} states it as Conjecture~5.1 in the same
words, and the author's later copy \cite[p.~11]{HagW} leaves the conjecture and
its lead-in unchanged. We use two statements about the pencil for a given $k$:
\begin{itemize}
\item[(D)] along every branch of zeros of $F_t$ in the open upper half-plane,
$\im z$ does not increase with $t$;
\item[(R)] a real zero of $F_t$ stays real as $t$ increases.
\end{itemize}
Statement (D) is the conjecture; (R) is the last clause of the lead-in. A
conjugate pair that leaves the real axis at some $t_0$ violates (R), and also
(D) just after $t_0$, since the imaginary part of its upper member rises from
$0$.

\subsection{What is shown here}\label{sec:shown}
Each item is labeled by its kind.
\begin{itemize}
\item \emph{Proved.} The pencil is $\Xi$ minus a level: $F_t = \Xi - L_t$ with
$L_t = u\Phi_{k+1} + Q_{k+1}$ positive on the real axis and
$\partial_t F_t = \Phi_{k+1} > 0$ there (Lemma~\ref{lem:level}). On the real
axis $F_t$ increases with $t$; its real zeros lie in the positive lobes of
$\Xi$, an even number in each and an odd number in the central one; and (R)
holds for $k$ if and only if every critical point of
$S_k := \Xi_{k+1}/\Phi_{k+1}$ on the real axis with value in $(0,1]$ is a
non-degenerate maximum (Theorem~\ref{thm:real}).
\item \emph{Proved under the hypothesis that every zero of $\Xi$ is real.} A
conjugate pair can leave the real axis only where the level satisfies an
explicit inequality (Proposition~\ref{prop:liftoff}). If the zeros of $\Xi$ are
also simple and the level is frozen at a constant that decreases, every
non-real zero descends along its own curve and lands at the top of its own lobe
(Proposition~\ref{prop:frozen}); this is known in substance \cite{CS}.
\item \emph{Proved.} If $\Xi$ had a simple zero $\beta$ with $\im\beta > 0$ and
$\im\Xi'(\beta) > 0$, then for every large $k$ the pencil would have a non-real
zero whose imaginary part strictly increases on the whole interval
$0 \le t \le 1$ (Proposition~\ref{prop:offaxis}).
\item \emph{Proved.} Along the pencil the coefficient of $1/x^2$ on the real
axis is negative and decreases in size (Lemma~\ref{lem:ct}).
\item \emph{Numerical, floating point, not interval arithmetic.} A census of the
branches of zeros of the pencil by two independent programs, and a test of the
real-axis criterion (Section~\ref{sec:census}).
\item \emph{Heuristic.} Why the true, slowly varying level does not change the
frozen picture, in three regimes (Section~\ref{sec:heur}).
\end{itemize}
Neither (D) nor (R) is proved here for any $k$; Baccaro's preprint
\cite{Bac} treats $k = 1$ (Section~\ref{sec:prior}).

%% SECTION 2
\section{The level form and the real axis}\label{sec:real}

\begin{lemma}\label{lem:phipos}
For $n \ge 2$ and real $x$, $\Phi_n(x) > 0$.
\end{lemma}

\begin{proof}
This is the P\'olya argument of \cite[Theorem~3.1]{Tya} applied to the single
kernel $\tphi_n$. For $x \ne 0$,
$\int_0^\infty \tphi_n(v)\cos(xv)\,dv
= x^{-1}\int_0^\infty(-\tphi_n'(v))\sin(xv)\,dv$ with $-\tphi_n'$ positive and
strictly decreasing, so the pieces of the sine integral over consecutive
half-periods alternate in sign with strictly decreasing sizes, the first
positive. At $x = 0$ the integral is that of a positive function.
\end{proof}

\begin{lemma}[the pencil is $\Xi$ minus a positive level]\label{lem:level}
For $k \ge 1$ and all $z$,
\[
F_t(z) = \Xi(z) - L_t(z), \qquad
L_t := u\Phi_{k+1} + Q_{k+1} = 2\int_0^\infty \kappa_t(v)\cos(zv)\,dv, \qquad
\kappa_t := u\tphi_{k+1} + \sum_{n > k+1}\tphi_n .
\]
For $0 \le u \le 1$ the kernel $\kappa_t$ is positive, strictly decreasing and
strictly convex, so $L_t(x) > 0$ for real $x$; and
$\partial F_t/\partial t = \Phi_{k+1}$, which is positive on the real axis.
\end{lemma}

\begin{proof}
$\Xi_k = \Xi - Q_k$ and $Q_k = \Phi_{k+1} + Q_{k+1}$. The kernel properties hold
termwise, since all indices are at least $2$; then the argument of
Lemma~\ref{lem:phipos} applies to $\kappa_t$, and Lemma~\ref{lem:phipos} itself
gives $\Phi_{k+1} > 0$.
\end{proof}

\begin{lemma}[the quotient form]\label{lem:quot}
Let $S_k := \Xi_{k+1}/\Phi_{k+1}$, a meromorphic function that is real on the
real axis, where $\Phi_{k+1} > 0$. Then $F_t(z) = 0$ if and only if
$S_k(z) = u$, at every $z$ with $\Phi_{k+1}(z) \ne 0$; at a common zero of
$\Xi_{k+1}$ and $\Phi_{k+1}$ the pencil vanishes for every $t$. Along a branch
of simple zeros $z(t)$ with $\Phi_{k+1}(z) \ne 0$,
\[
\frac{dz}{dt} = -\frac{\Phi_{k+1}(z)}{\partial_z F_t(z)} = -\frac{1}{S_k'(z)},
\qquad\text{so}\qquad
\frac{d(\im z)}{dt} = \frac{\im S_k'(z)}{|S_k'(z)|^2}.
\]
Hence (D) for $k$ is the statement that $\im S_k'(z) \le 0$ at every $z$ with
$\im z > 0$ and $S_k(z) \in (0,1)$, together with the behavior of the branches
at multiple zeros.
\end{lemma}

\begin{proof}
$F_t = \Xi_{k+1} - u\Phi_{k+1}$; differentiate $S_k(z(t)) = 1 - t$.
\end{proof}

Let $0 < \gamma_1 < \gamma_2 < \cdots$ be the positive real zeros of $\Xi$; a
\emph{positive lobe} is an interval $(\gamma_j, \gamma_{j+1})$ on which
$\Xi > 0$, and the \emph{central lobe} is $(0, \gamma_1)$.

\begin{theorem}\label{thm:real}
Let $k \ge 1$ and $0 \le t \le 1$.
\begin{itemize}
\item[(a)] For real $x$, $F_t(x) = \Xi(x) - L_t(x)$ with $L_t(x) > 0$, and
$t \mapsto F_t(x)$ is strictly increasing.
\item[(b)] Every real zero of $F_t$ lies in an open interval on which $\Xi > 0$.
Counted with multiplicity, each positive lobe $(\gamma_j, \gamma_{j+1})$
contains an even number of zeros of $F_t$ and the central lobe an odd number;
the number of zeros of $F_t$ in $(0,\infty)$ is finite and odd.
\item[(c)] The sets $P_t := \{x \in \R : F_t(x) > 0\}$ increase with $t$:
$P_s \subset P_t$ for $s < t$, and every boundary point of $P_s$ lies in the
interior of $P_t$.
\item[(d)] Let $x_0$ be a real zero of $F_{t_0}$ of multiplicity $m$, with
$F_{t_0}(x) = c(x - x_0)^m(1 + o(1))$. If $m$ is odd, $F_t$ has an odd number
of real zeros near $x_0$ for $t$ near $t_0$, counted with multiplicity; for
$t \ne t_0$ the $m$ zeros near $x_0$ are
$x_0 + \bigl(-(t - t_0)\Phi_{k+1}(x_0)/c\bigr)^{1/m}(1 + o(1))$ (the $m$
values of the root), exactly one of them real and $m - 1$ non-real, so for odd
$m \ge 3$ a conjugate pair lands at $t_0$ and leaves the axis at once. The same
expansion gives, for even $m \ge 4$, two real and $m - 2$ non-real zeros on the
side of $t_0$ where $F_t(x_0)$ has the sign opposite to $c$, and $m$ non-real
zeros on the other side. If $m = 2$ and $F_{t_0} \le 0$ near $x_0$ (a local
maximum with value $0$), then for $t$ slightly above $t_0$ there are two simple
real zeros near $x_0$ and for $t$ slightly below there are none: a conjugate
pair arrives, a \emph{landing}. If $m = 2$ and $F_{t_0} \ge 0$ near $x_0$ (a
local minimum with value $0$), then for $t$ slightly below $t_0$ there are two
real zeros near $x_0$ and for $t$ slightly above there are none: a conjugate
pair leaves the axis, a \emph{lift-off}.
\item[(e)] Hence (R) holds for $k$ if and only if, for every $0 \le t < 1$,
every multiple real zero of $F_t$ is a double zero at which $F_t$ is locally
$\le 0$; equivalently, every critical point of $x \mapsto S_k(x)$ with value in
$(0,1]$ is a non-degenerate local maximum ($S_k'' < 0$). The weaker conditions
``no $F_t$, $0 \le t < 1$, has a real zero of even multiplicity at which it is
locally non-negative'', ``no two components of $P_t$ merge as $t$ increases''
and ``$S_k$ has no local minimum with value in $(0,1]$'' are equivalent to one
another and necessary for (R); they are sufficient when no $F_t$,
$0 \le t < 1$, has a real zero of multiplicity at least $3$ (no critical point
of $S_k$ with $S_k'' = 0$ and value in $(0,1]$), since at such a zero at least
one conjugate pair leaves the axis just after $t_0$, by (d).
\end{itemize}
\end{theorem}

\begin{proof}
(a) Lemma~\ref{lem:level}.
(b) Where $\Xi \le 0$, $F_t = \Xi - L_t < 0$; in particular $F_t(\gamma_j) < 0$
for every $j$. At the origin $F_t(0) \ge \Xi_1(0) > 0$, because
$F_t - \Xi_1 = \sum_{2 \le n \le k}\Phi_n + t\Phi_{k+1} \ge 0$ at real points
(Lemma~\ref{lem:phipos}) and $\Xi_1(0) = \Xi(0) - Q_1(0) \ge 0.4971 - c_1 > 0$,
where $c_1 = \sum_{n \ge 2}(4\pi n^2 - 1)e^{-\pi n^2} \ge Q_1(0)$
\cite[proof of Corollary~3.5]{Tya}. For large $x$, $x^2F_t(x)$ tends to a
negative limit: $\Xi(x)$ decays exponentially, and two integrations by parts,
as in the proof of \cite[Proposition~3.4]{Tya} applied to the kernel
$\kappa_t$, give $\lim_{x\to\infty} x^2L_t(x) = -2\kappa_t'(0) > 0$. So the
zeros of $F_t$ in $[0,\infty)$ lie in a bounded set and are finite in number.
The parity statements follow from the signs at the ends of each lobe.
(c) $F_s(x) > 0$ implies $F_t(x) > F_s(x) > 0$; a boundary point of $P_s$ has
$F_s = 0$ there, so $F_t > 0$ there.
(d) $F_t(x) = F_{t_0}(x) + (t - t_0)\Phi_{k+1}(x)$ with $\Phi_{k+1}(x_0) > 0$;
near $x_0$ this is
$c(x - x_0)^m(1 + o(1)) + (t - t_0)\Phi_{k+1}(x_0)(1 + o(1))$, and the cases
are read off; Weierstrass preparation gives exactly $m$ zeros near $x_0$, real
or in conjugate pairs.
(e) From (d), since $F_t(x) = \Phi_{k+1}(x)\bigl(S_k(x) - u\bigr)$ on the real
axis.
\end{proof}

\begin{remark}\label{rem:real}
Zeros of multiplicity at least $3$ are covered by the same expansion, and at
every such zero (R) and (D) fail, by (d); none occurs in any computation of
Section~\ref{sec:census}, and a degenerate critical point of $S_k$
($S_k' = S_k'' = 0$ at one $x$) is non-generic. Theorem~\ref{thm:real} does not
decide (R): it reduces it to ``every critical point of $S_k$ with value in
$(0,1]$ is a non-degenerate maximum''. A scan that finds no local minimum with
value in $(0,1]$ checks this up to degenerate critical points, which it should
flag ($|S_k''|$ small at an extremum).
\end{remark}

\begin{proposition}[conditional on the zeros of $\Xi$ being real]\label{prop:liftoff}
Assume every zero of $\Xi$ is real. If $F_{t_0}$ has a lift-off at $x_0$
(Theorem~\ref{thm:real}(d)), then
\[
-(\log L_{t_0})''(x_0) \ \ge\ \sum_{\gamma > 0}
\Bigl[(x_0 - \gamma)^{-2} + (x_0 + \gamma)^{-2}\Bigr].
\]
\end{proposition}

\begin{proof}
At a lift-off $\Xi = L$, $\Xi' = L'$ and $\Xi'' \ge L''$ at $x_0$, with
$L = L_{t_0} > 0$; so
$(\log\Xi)''(x_0) = \Xi''/\Xi - (\Xi'/\Xi)^2 \ge L''/L - (L'/L)^2
= (\log L)''(x_0)$. With only real zeros,
$\Xi(z) = \Xi(0)\prod_{\gamma > 0}(1 - z^2/\gamma^2)$ ($\Xi$ is even, of
order $1$, and $\sum\gamma^{-2}$ converges), so
$(\log\Xi)''(x) = -\sum_{\gamma>0}\bigl[(x - \gamma)^{-2} + (x + \gamma)^{-2}\bigr]$.
\end{proof}

So, under that hypothesis, (R) follows from the inequality
$-(\log L_t)''(x) < \sum_{\gamma > 0}[(x - \gamma)^{-2} + (x + \gamma)^{-2}]$
at the real points where $\Xi = L_t$. The right side is at least
$2\sum\gamma^{-2} = 0.0462\ldots$ in the central lobe and at least
$8(\gamma_{j+1} - \gamma_j)^{-2}$ in the lobe $(\gamma_j, \gamma_{j+1})$ (the
two endpoint terms, smallest at the midpoint). The left side is,
heuristically, of the order of the squared width of the kernel $\kappa_t$,
about $1/(2\pi^2(k+1)^4)$. A proved bound for $-(\log L_t)''$ is not given
here; with the hypothesis of Proposition~\ref{prop:liftoff} it would give (R)
for a range of $k$. Whether a hypothesis on the zeros of $\Xi$ up to a finite
height suffices is not shown here: the product formula would need a remainder
for the zeros above that height. The inequality behind
Proposition~\ref{prop:liftoff} is tabulated numerically in
Section~\ref{sec:census}.

%% SECTION 3
\section{The frozen level}\label{sec:frozen}

By Lemma~\ref{lem:level} the zeros of $F_t$ are the solutions of
$\Xi(z) = L_t(z)$. If the level were a constant $c > 0$ that decreases with
$t$, (D) and (R) would be statements about the solutions of $\Xi = c$. When the
zeros of $\Xi$ are real these are known in substance, from the formula for
$\im f'/f$ in Titchmarsh \cite[\S8.52, p.~266]{Tit} and from the level-curve
statements of Csordas and Smith \cite[(2.5), p.~604; (3.3)--(3.4), p.~609]{CS};
we state them in the form needed here, with a short proof.

\begin{lemma}[Titchmarsh {\cite[\S8.52, p.~266]{Tit}}]\label{lem:tit}
Let $f$ be a real entire function of the form
$f(z) = Cz^m e^{-az^2 + bz}\prod_j(1 - z/a_j)e^{z/a_j}$ with $a \ge 0$, $b$ and
the $a_j$ real, $\sum_j a_j^{-2} < \infty$, and at least one zero or $a > 0$.
Then $\im f'(z)/f(z) < 0$ for $\im z > 0$. If $c \ne 0$ is real and $z(c)$ is
a branch of simple solutions of $f(z) = c$ in the open upper half-plane, then
$d(\im z)/d|c| > 0$: as $|c|$ decreases, the solution moves toward the real
axis.
\end{lemma}

\begin{proof}
$f'/f = m/z - 2az + b + \sum_j[1/(z - a_j) + 1/a_j]$, and each of the terms
$m/z$, $-2az$, $1/(z - a_j)$ has non-positive imaginary part for $\im z > 0$,
with strict inequality for at least one. Then $dz/dc = 1/f'(z) = (1/c)(f/f')(z)$
because $f(z) = c$, and $\im(f/f') = -\im(f'/f)/|f'/f|^2 > 0$.
\end{proof}

\begin{proposition}[the solutions of $\Xi = c$ when the zeros of $\Xi$ are real; in substance {\cite{CS}}]\label{prop:frozen}
Assume every zero of $\Xi$ is real and simple, and let $M_m$ be the maximum of
$\Xi$ over the $m$-th positive lobe $(\gamma_{2m}, \gamma_{2m+1})$, $m \ge 1$.
\begin{itemize}
\item[(i)] For every $y > 0$ the continuous argument of $\Xi$ along the
horizontal line through $iy$, normalized by $\arg\Xi(iy) = 0$, is strictly
decreasing in $x$ and tends to $-\infty$; so for each $m \ge 1$ there is exactly
one point $x_m(y) + iy$ with $x_m(y) > 0$ and $\arg\Xi = -2\pi m$. The curves
$C_m := \{x_m(y) + iy : y > 0\}$ are disjoint real-analytic graphs over the
$y$-axis, and
$\{z : \re z > 0,\ \im z > 0,\ \Xi(z) > 0\} = \bigcup_{m \ge 1} C_m$.
\item[(ii)] Along $C_m$ the positive number $\Xi$ is strictly increasing in
$y$; as $y \to 0$ the curve tends to the critical point of $\Xi$ in the $m$-th
positive lobe, and $\Xi \to M_m$.
\item[(iii)] Hence, for a constant $c > 0$, the zeros of $\Xi - c$ in the open
first quadrant are in bijection with the positive lobes with $M_m < c$: one
zero on each such $C_m$, none on the others, which carry two real zeros of
$\Xi - c$ in the lobe when $M_m > c$. As $c$ decreases each such zero moves
down its own curve $C_m$, strictly monotonically in $\im z$, and reaches the
real axis exactly when $c = M_m$, at the critical point of the lobe.
\end{itemize}
\end{proposition}

\begin{proof}
(i) Under the hypothesis $\Xi(z) = \Xi(0)\prod(1 - z^2/\gamma^2)$, and
$\partial_x\arg\Xi(x + iy) = \im\Xi'/\Xi \le -\sum_\gamma y/((x - \gamma)^2 + y^2)$
by Lemma~\ref{lem:tit}; the integral over $x$ diverges because there are
infinitely many $\gamma$. Also
$\Xi(iy) = 2\int_0^\infty\sum_n\tphi_n(v)\cosh(yv)\,dv > 0$. The implicit
function theorem gives the graphs.
(ii) $\Xi' \ne 0$ off the real axis (Lemma~\ref{lem:tit}), so $\Xi$, real on
$C_m$, is strictly monotone along it, and by Lemma~\ref{lem:tit} it increases
with $y$ (that such level curves have no horizontal tangent is
\cite[(3.3)--(3.4)]{CS}; the direction is Lemma~\ref{lem:tit}).
Near a real point that is neither a critical point nor a zero of $\Xi$, and
near a simple zero, $\Xi$ is injective and maps the real axis onto real values;
so $C_m$ can approach the real axis only at a critical point where $\Xi > 0$.
As $y \to 0$ the argument tends, locally uniformly away from the zeros, to its
boundary values ($-\pi$ per simple zero passed: $-2\pi m$ on the $m$-th positive
lobe, $-2\pi m \pm \pi$ on the adjacent negative lobes), so every limit point of
$C_m$ is a critical point of $\Xi$ in the closed lobe, not an endpoint. Under
the hypothesis
$(\Xi'/\Xi)'(x) = -\sum_\gamma[(x - \gamma)^{-2} + (x + \gamma)^{-2}] < 0$ on
each lobe, so $\Xi'/\Xi$ falls strictly from $+\infty$ to $-\infty$ there: each
lobe contains exactly one critical point, where $\Xi'' = \Xi\cdot(\Xi'/\Xi)' < 0$
(a non-degenerate maximum, value $M_m$). Hence $C_m$ tends to it.
(iii) For $y \ge 2\pi e^2$ and $\sigma := \frac12 + y$ one has
$\Xi(x + iy) = \overline{\xi(\sigma + ix)}$ and, for $x \ge 0$,
$\arg\xi(\sigma + ix) \ge (x/2)\log(\sigma/(2\pi e)) - \pi/2$, from
$\arg\Gamma(s/2) = \frac12\int_0^x \re\psi((\sigma + i\tau)/2)\,d\tau$ with
$\re\psi(a + ib) \ge \psi(a) \ge \log a - 1/a$, $\arg s(s-1) \in [0,\pi)$, and
$|\arg\zeta(s)| < \pi/2$ for $\sigma \ge 2$. So $x_m(y) \le 4\pi m + \pi$, and
$|\xi(\sigma + ix)| \to \infty$ uniformly for $0 \le x \le 4\pi m + \pi$ as
$\sigma \to \infty$. Thus $\Xi$ maps $C_m$ increasingly onto $(M_m, \infty)$,
which gives one zero on each such $C_m$; the two real zeros in a lobe with
$M_m > c$ come from its single critical point.
\end{proof}

\emph{What this says about the pencil.} If $L_t$ were a constant that
decreases with $t$, (D) and (R) would be exactly
Proposition~\ref{prop:frozen}(iii): every non-real zero descends, each lands at
the top of its own lobe when the level passes that top, and landed zeros stay
real. The order in which the lobes are reached is the order of their heights
$M_m$, not of their positions. This is how Haglund's Conjecture~1 (for every
$N$, the zeros of $\Xi_N$ in the closed first quadrant, listed by increasing
real part, have nondecreasing imaginary parts \cite[p.~3]{Hag}) fails at
$N = 27$ \cite{Tya}: a low lobe to the left of a higher one. It is compatible
with (D): the zero over the low lobe is still descending when
the zeros of the higher lobe to its right have landed. The true level $L_t(z)$
is not constant: on the real axis it varies on the scale $2\pi(k+1)^2$, and off
the axis it has a small phase (Section~\ref{sec:heur}).

\emph{The exact criterion, with no hypothesis.} At a simple zero $z$ of $F_t$
with $\im z > 0$ that is not a common zero of $\Xi$ and $L_t$,
$\Xi(z) = L_t(z) \ne 0$ and
\[
\frac{dz}{dt} = -\frac{\rho(z)}{\Lambda(z)}, \qquad
\rho := \frac{\Phi_{k+1}}{L_t}, \qquad
\Lambda := \frac{\Xi'}{\Xi} - \frac{L_t'}{L_t},
\]
so the zero descends if and only if $\im(\rho/\Lambda) > 0$. In the frozen
model $\rho > 0$ and $\Lambda = \Xi'/\Xi$.

%% SECTION 4
\section{A zero of \texorpdfstring{$\Xi$}{Xi} off the real axis}\label{sec:offaxis}

\begin{lemma}[$\Phi_n$ is nearly constant on a fixed disc]\label{lem:disc}
Let $R > 0$ and $n \ge 2$ with $\pi n^2 \ge R + \frac12$. Then for $|z| \le R$,
\[
\left|\frac{\Phi_n(z)}{\Phi_n(0)} - 1\right| \le \frac{3.1\,R^2}{\pi^2 n^4}.
\]
\end{lemma}

\begin{proof}
Put $X = \pi n^2$. For $v \ge 0$,
$|\cos(zv) - 1| \le \cosh(Rv) - 1 \le \frac12 R^2v^2e^{Rv}$. With
$y = Xe^{2v}$ (so $v = \frac12\log(y/X)$, $dv = dy/2y$) one has
$\tphi_n(v)\,dv = (2y - 3)(y/X)^{1/4}e^{-y}\,dy$. Hence
$\int_0^\infty\tphi_n\,dv \ge \int_X^\infty(2y - 3)e^{-y}\,dy = (2X - 1)e^{-X}$.
With $w = y - X$: $\log(y/X) \le w/X$ and
$(y/X)^{1/4 + R/2} \le e^{(1/4 + R/2)w/X} \le e^{w/2}$ because
$X \ge \frac12 + R$; so
\[
\int_0^\infty\tphi_n(v)\cdot\tfrac12R^2v^2e^{Rv}\,dv
\le \frac{R^2}{8X^2}\,e^{-X}\int_0^\infty(2X + 2w)\,w^2e^{-w/2}\,dw
= \frac{R^2}{8X^2}(32X + 192)\,e^{-X}.
\]
The quotient is at most
$(R^2/8X^2)(32X + 192)/(2X - 1) \le 3.1R^2/X^2$ for $X \ge 4\pi$.
\end{proof}

\emph{Consequences.} For a fixed disc $|z| \le R$ and $k \to \infty$, uniformly
in $0 \le t \le 1$: $\Phi_{k+1}(z)/\Phi_{k+1}(0) \to 1$;
$\sup|L_t| \le \sum_{n > k}\sup|\Phi_n| \le \Phi_{k+1}(0)(1 + o(1)) \to 0$,
because the ratios $\Phi_{n+1}(0)/\Phi_n(0)$ tend to $0$, since
$(2X - 1)e^{-X} \le \frac12\Phi_n(0) \le (2X + 3)e^{-X}$ for
$X = \pi n^2 \ge 4\pi$ (the upper bound from $2y - 3 \le 2y$ and
$(y/X)^{1/4} \le e^{w/4X}$, which give
$e^{-X}[2X/(1 - \varepsilon) + 2/(1 - \varepsilon)^2]$ with
$\varepsilon = 1/4X \le 0.02$); and $L_t' \to 0$ on any smaller disc, by
Cauchy's estimate.

\begin{proposition}\label{prop:offaxis}
Let $\beta$ be a simple zero of $\Xi$ with $\im\beta > 0$ and
$\im\Xi'(\beta) \ne 0$. Then there are a disc $D$ about $\beta$ and $k_0$ such
that for every $k \ge k_0$ and every $t \in [0,1]$ the function
$F_t = \Xi_k + t\Phi_{k+1}$ has exactly one zero $z_k(t)$ in $D$, it is
simple, $t \mapsto z_k(t)$ is continuously differentiable, and
\[
\operatorname{sign}\frac{d(\im z_k(t))}{dt} = \operatorname{sign}\im\Xi'(\beta)
\qquad\text{for all } t \in [0,1].
\]
In particular, if $\Xi$ had a simple zero $\beta$ in the open upper half-plane
with $\im\Xi'(\beta) > 0$, then for every $k \ge k_0$ the pencil would have a
non-real zero whose imaginary part strictly increases on the whole interval
$0 \le t \le 1$, and (D) would fail for all large $k$.
\end{proposition}

\begin{proof}
Choose $D$ with closure in the open upper half-plane, containing no other zero
of $\Xi$ and no zero of $\Xi'$, and let $\delta = \min_{\partial D}|\Xi| > 0$.
By the consequences of Lemma~\ref{lem:disc}, $\sup_{\overline D}|L_t| < \delta$
for $k \ge k_1$ and all $t$, so by Rouch\'e's theorem $F_t = \Xi - L_t$ has
exactly one zero $z_k(t)$ in $D$, and it is simple; it depends continuously
differentiably on $t$ by the implicit function theorem, with
$dz_k/dt = -\Phi_{k+1}(z_k)/(\Xi'(z_k) - L_t'(z_k))$. As $k \to \infty$,
uniformly in $t$: $z_k(t) \to \beta$ (the zero of $\Xi - L_t$ in $D$ tends to
the zero of $\Xi$ as $\sup|L_t| \to 0$),
$\Xi'(z_k) - L_t'(z_k) \to \Xi'(\beta) \ne 0$ and
$\Phi_{k+1}(z_k)/\Phi_{k+1}(0) \to 1$. Therefore
\[
\frac{1}{\Phi_{k+1}(0)}\,\frac{dz_k}{dt} \to -\frac{1}{\Xi'(\beta)}
\qquad\text{uniformly on } [0,1],
\]
and $\im(-1/\Xi'(\beta)) = \im\Xi'(\beta)/|\Xi'(\beta)|^2 \ne 0$ has the sign
of $\im\Xi'(\beta)$.
\end{proof}

\begin{remark}\label{rem:offaxis}
(1) The mirror pencil $\Xi + L_t$ (the weights $1, \ldots, 1, 2 - t, 2, 2,
\ldots$ on $\Phi_1, \Phi_2, \ldots$, still a series of Haglund's form
\cite[(53)]{Hag} with non-negative coefficients) has
$dz/dt = +\Phi_{k+1}/(\Xi' + L_t')$, so the sign is reversed: a simple off-axis
zero with $\im\Xi'(\beta) \ne 0$ produces an ascending branch, for all large
$k$, in exactly one of the two pencils.
(2) Not covered: $\im\Xi'(\beta) = 0$, and multiple zeros off the axis. For a
zero of multiplicity $m \ge 2$ the $m$ nearby zeros of $F_t$ sit at $\beta$
plus the $m$-th roots of $L_t(\beta)/c$, $c = \Xi^{(m)}(\beta)/m!$, and approach
$\beta$ radially as $t$ increases ($|L_t(\beta)|$ decreases with $u$) and as
$k$ grows; at least one approaches from below, and so ascends in $t$, when
$m \ge 3$, and when $m = 2$ unless $c$ is a positive real number. This is a
sketch, not a proof.
(3) Proposition~\ref{prop:offaxis} concerns large $k$ only; it says nothing
about a fixed $k$.
(4) By the reflection $\Xi(-\bar z) = \overline{\Xi(z)}$, the zero $-\bar\beta$
has $\Xi'(-\bar\beta) = -\overline{\Xi'(\beta)}$, with the same sign of the
imaginary part; the two mirror zeros behave alike.
\end{remark}

Read in the other direction, Proposition~\ref{prop:offaxis} says that (D) for
all large $k$ excludes one class of zeros of $\Xi$ off the real axis: the
simple ones with $\im\Xi'(\beta) > 0$. Simple off-axis zeros with
$\im\Xi'(\beta) < 0$ descend in this pencil and are not excluded, and multiple
ones are not treated. In the model where the level is a constant, real and
simple zeros of $\Xi$ give (D) and (R) (Proposition~\ref{prop:frozen}); that is
a model, not a statement about the pencil (Section~\ref{sec:heur}).

%% SECTION 5
\section{The census}\label{sec:census}

Everything in this section is numerical: floating point, not interval
arithmetic. Two programs were written independently.
\emph{Program~A} evaluates $\Phi_n$ and $S_k$ in Arb ball arithmetic
\cite{Arb, FLINT}, using every value only when its relative radius is below
$2^{-40}$; it follows each branch as the level curve $\im S_k = 0$
(Lemma~\ref{lem:quot}) from a zero of $\Xi_k$ to a landing or to a zero of
$\Xi_{k+1}$, and counts zeros by the argument principle in floating point.
\emph{Program~B} uses mpmath \cite{mpmath} only, written from Haglund's
definitions without sight of program~A; it follows each branch by Newton
continuation on a grid in $\tau = -\log(1 - t)$, and reproduces Haglund's
$19$ printed zeros of $\Xi_1$ to $7\cdot10^{-22}$, the largest real zeros of
$\Xi_1, \ldots, \Xi_4$ in his table, and the certified zeros of $\Xi_{27}$ of
\cite{Tya}.

\begin{table}[ht]
\centering\small
\setlength{\tabcolsep}{3.6pt}
\begin{tabular}{rllrrrrrrrrl}
\toprule
$k$ & $\re z$ & $\im z$ & real & non-real & B & L & E & X & $\Delta^+$ & $m$ & prog.\\
\midrule
1 & $[0, 86.55]$ & $[0, 41]$ & 1/7 & 15/11 & 15 & 3 & 11 & 1 & $-3.7\cdot10^{-7}$ & 0.708 & A, B \\
2 & $[0, 130.53]$ & $[0, 50]$ & 7/15 & 23/18 & 23 & 4 & 18 & 1 & $-3.1\cdot10^{-7}$ & 0.806 & A, B \\
3 & $[0, 187.08]$ & $[0, 59]$ & 15/31 & 35/25 & 35 & 8 & 25 & 2 & $-3.3\cdot10^{-7}$ & 0.812 & A, B \\
4 & $[0, 256.19]$ & $[0, 68]$ & 31/53 & 47/34 & 47 & 11 & 34 & 2 & $-3.2\cdot10^{-7}$ & 0.797 & A \\
5 & $[0, 337.88]$ & $[0, 79]$ & 53/79 & 63/47 & 63 & 13 & 47 & 3 & $-3.0\cdot10^{-7}$ & 0.828 & A \\
6 & $[0, 432.12]$ & $[0, 90]$ & 79/113 & 82/62 & 82 & 17 & 62 & 3 & $-3.0\cdot10^{-7}$ & 0.832 & A \\
7 & $[0, 538.94]$ & $[0, 102]$ & 113/155 & 103/80 & 103 & 21 & 80 & 2 & $-3.1\cdot10^{-7}$ & 0.842 & A \\
8 & $[0, 658.32]$ & $[0, 115]$ & 155/207 & 127/98 & 127 & 26 & 98 & 3 & $-3.0\cdot10^{-7}$ & 0.850 & A \\
\midrule
15 & $[984, 1196]$ & $[0, 55]$ & 37/145$^*$ & 74/16 & 74 & 54 & 16 & 4 & $-3.0\cdot10^{-7}$ & 0.880 & A \\
20 & $[1724, 1976]$ & $[0, 61]$ & 40/196$^*$ & 101/17 & 101 & 78 & 17 & 6 & $-3.0\cdot10^{-7}$ & 0.876 & A \\
26 & $[2876, 3176]$ & $[0, 69]$ & 45/263$^*$ & 133/17 & 133 & 109 & 17 & 7 & $-3.0\cdot10^{-7}$ & 0.868 & A \\
27 & $[3096, 3404]$ & $[0, 70]$ & 48/276$^*$ & 137/16 & 137 & 114 & 16 & 7 & $-3.0\cdot10^{-7}$ & 0.851 & A \\
35 & $[5144, 5516]$ & $[0, 80]$ & 53/363$^*$ & 182/19 & 182 & 155 & 19 & 8 & $-3.0\cdot10^{-7}$ & 0.876 & A \\
50 & $[10364, 10856]$ & $[0, 99]$ & 58/546$^*$ & 274/19 & 274 & 244 & 19 & 11 & $-3.0\cdot10^{-7}$ & 0.857 & A \\
\bottomrule
\end{tabular}
\caption{Branches of zeros of $F_t = \Xi_k + t\Phi_{k+1}$, $0 \le t \le 1$, in
the window $W$ given by the ranges of $\re z$ and $\im z$. Upper block: complete
windows, $\re z \le 2\pi(k+2)^2 + 30$. Lower block: windows
$[4(k+1)^2 - 40,\ 4(k+2)^2 + 40]$ around the landings, branches that start in
them. Real and non-real: numbers of zeros of $\Xi_k$ / $\Xi_{k+1}$ in $W$
(${}^*$: real zeros in the window). B: branches followed (the non-real zeros of
$\Xi_k$ in $W$); L: landings; E: ends at a non-real zero of $\Xi_{k+1}$;
X: exits through the side of $W$. $\Delta^+$: the largest change of $\im z$
over one step, over all steps of all branches (negative: $\im z$ decreased at
every step). $m$: the smallest margin $-\im S_k'/|S_k'|$, the sine of the angle
of descent, above height $10^{-6}$. Floating point, not interval arithmetic.}
\label{tab:census}
\end{table}

\emph{Readings} (each a statement about its window). In every row of
Table~\ref{tab:census} the imaginary part decreased at every step of every
branch, and the margin stays above $0.7$; the number of landings equals half the
difference of the real counts; and every non-real zero of $\Xi_{k+1}$ in the
window is the end of a followed branch (program~A's completeness checks; for
$k = 6$ one branch reaches $t = 1$ at the zero $432.1548 + 63.0618i$ of $\Xi_7$,
just beyond the window, and is counted as an exit). Where both programs ran
($k = 1, 2, 3$, program~B with window heights $45$, $45$, $60$) they agree on
every count, and on all $15$ landings to the digits program~B printed:
$|\Delta x^*| \le 3.9\cdot10^{-7}$, $|\Delta\tau^*| \le 4.8\cdot10^{-6}$.
Program~B also finds $\im(dz/dt) < 0$ at all $3{,}605$, $6{,}886$ and
$11{,}974$ grid values of its branches for $k = 1, 2, 3$. The real counts in
the upper block are those of Haglund's table \cite[p.~4]{HagW}, which has $31$
at $N = 4$; the arXiv and journal versions print $32$ there
\cite[p.~4]{Hag}.

\emph{The real axis.} For $k = 1, \ldots, 50$, on $[0, 4(k+2)^2 + 40]$, program~A
finds that $S_k$ has no local minimum with value in $(0,1)$ and that its local
maxima with value in $(0,1)$ are $(R_{k+1} - R_k)/2$ in number, where $R_N$ is
the number of real zeros of $\Xi_N$ found on that interval (grid of $1/20$ of
the mean zero spacing $2\pi/\log(x/2\pi)$, the part beyond $4(k+1)^2 - 40$
scanned again at $1/40$). Each $R_N$ is found twice, on the intervals of the
pencils $k = N - 1$ and $k = N$, with the same value; for $N = 1, \ldots, 51$:
1, 7, 15, 31, 53, 79, 113, 155, 207, 263, 327, 401, 483, 575, 673, 781, 899,
1027, 1163, 1307, 1463, 1631, 1807, 1991, 2189, 2393, 2611, 2839, 3075, 3323,
3581, 3851, 4133, 4423, 4727, 5037, 5363, 5699, 6045, 6405, 6771, 7157, 7549,
7951, 8369, 8797, 9237, 9689, 10155, 10629, 11117 --- all odd, as
\cite[Corollary~3.5]{Tya} requires. Program~B finds the same for $k = 1, 2, 3$
on $[0, 98.55]$, $[0, 142.5]$, $[0, 199.1]$. A scan for minima checks the
criterion of Theorem~\ref{thm:real}(e) up to degenerate critical points
(Remark~\ref{rem:real}). The inequality in the first step of the proof of
Proposition~\ref{prop:liftoff}, which uses no hypothesis (a lift-off at $x_0$
forces $(\log\Xi)''(x_0) \ge (\log L_t)''(x_0)$), was tabulated on the part of
$[0, 4(k+2)^2 + 40]$ where $\Xi > 0$ (step $0.02$, $u \in \{0, \frac14,
\frac12, \frac34, 1\}$): there $-(\log L_u)''$ is at most $0.066$ times
$-(\log\Xi)''$ at $k = 1$ and at most $5.6\cdot10^{-6}$ times at $k = 20$.
